\documentclass[11pt]{article}
\usepackage[margin=1in]{geometry}
\usepackage{amsmath,amssymb,amsthm,mathtools}
\usepackage[hidelinks]{hyperref}
\newtheorem{theorem}{Theorem}
\newtheorem{lemma}[theorem]{Lemma}
\newtheorem{corollary}[theorem]{Corollary}
\title{Prime-power polygons and fold data for the tilted cubic moments}
\author{Internal proof and exact certificates; no priority claim}
\date{17 November 2026}
\begin{document}
\maketitle
The current ordinary-factorial cubic frontier is
\[
 T_n(t)=\sum_{j=0}^n\binom nj(n)_{2j}t^j,\quad n\ge1.
\]
The \href{DYADIC_BINOMIAL_SFC.tex}{lower-exponent-one identities} prove
that $X(A+BX^2)$ satisfies Strong Factorial at every start if and only
if every $T_n$ is squarefree. The following results close unbounded index
families and supply exact arithmetic and analytic structures; the full
simplicity assertion and unrestricted FC/SFC remain open.

\section{A uniform shifted polygon, with both residual types}
\begin{theorem}\label{thm:shift}
Let $P=p^r$, $r\ge1$, $a\ge1$, $i\ge0$, where $p$ is prime and
$p>3\max(a,i)$. Put $n=aP+i$ and $S_r=(P-1)/(p-1)$.
The Newton polygon of $T_n$ has a horizontal edge of length $i$ if
$i>0$, followed by an edge of width $aP$ and reduced slope
$2S_r/P$, whose denominator is $P$.
The horizontal residual is $T_i$ (with $T_0=1$).
The nonhorizontal residual, normalized to constant term one, is
\[
 \begin{cases}
 \displaystyle \sum_{v=0}^a\binom av(a)_{2v}Y^v=T_a(Y),&i=0,\\[3pt]
 \displaystyle \sum_{v=0}^a\binom av(a+1)_{2v}Y^v,&i>0.
 \end{cases}                                                   \tag{1}
\]
All residual statements are in $\mathbb F_p$; neither requires assuming
that the displayed residuals are squarefree.
\end{theorem}
\begin{proof}
The coefficient is
$A_j=\binom nj(n+2j-1)!/(n-1)!$.
All factorial arguments are below $pP$, since
$3aP+3i-1\le(p-1)P+p-2<pP$.
For a fixed level $b=p^h$, $1\le h\le r$, write $j=ub+c$,
$0\le c<b$. At this level the binomial valuation is
$\mathbf1_{\{c>i\}}$, because $n\bmod b=i$.
If $i\ge1$, the total level contribution is
\[
 2u+\left\lfloor\frac{i-1+2c}{b}\right\rfloor
                         +\mathbf1_{\{c>i\}}
                 \ \ge\ \frac{2(j-i)}b.                       \tag{2}
\]
For $c\le i$, the extra floor and indicator are zero and the right-hand
remainder is nonpositive. For $c>i$, if the floor is zero then
$2(c-i)<b$; if it is positive, the last two terms are at least two,
whereas $2(c-i)/b<2$.
If $i=0$, the contribution is $2u$ at $c=0$, and otherwise is
\[
             2u+2+\left\lfloor\frac{2c-1}b\right\rfloor
                              >\frac{2j}b.                    \tag{3}
\]
Summing (2) or (3) over $h$ proves
$v_p(A_j)\ge(2S_r/P)(j-i)$ for $j\ge i$.
The coefficients through degree $i$ are units when $i>0$, because
$j\le i<p$ and the short factorial block lies between residues $i$ and
$3i-1<p$. At every $j=i+vP$, $0\le v\le a$, the valuation is exactly
$2vS_r$. These points, including both endpoints, prove the polygon.
Its reduced denominator is $P$ since $p\ge5$ and $S_r\equiv1\pmod p$.

For completeness, the normalized unit part of $N!$ satisfies
\[
 p^{-v_p(N!)}N!\equiv(-1)^{v_p(N!)}\prod_\ell N_\ell!\pmod p,
                                                               \tag{4}
\]
where $N_\ell$ are its base-$p$ digits. Separate the multiples of $p$
and apply Wilson's theorem repeatedly to prove (4).
For $0<c<p$, $0\le h<p$, (4) gives units
$(-1)^{cS_r}c!h!$ at $N=cP+h$, and
$(-1)^{cS_r}(c-1)!$ at $N=cP-1$.
Also $\binom{aP+i}{vP+i}\equiv\binom av\pmod p$, by
$(1+X)^{aP+i}\equiv(1+X^P)^a(1+X)^i$.
At $i=0$ the normalized factorial ratio at $j=vP$ is therefore
$(a+2v-1)!/(a-1)!$.
At $i>0$, dividing by the coefficient at $j=i$ gives instead
$(a+2v)!/a!$; the common factor $(3i-1)!/(i-1)!$ cancels.
The signs in (4) cancel because the valuation difference is $2vS_r$.
These are exactly the two residuals in (1).
The horizontal coefficients reduce to
$\binom ij(i)_{2j}$; all higher coefficients have positive valuation.
\end{proof}

Every rational square divisor satisfies the occupancy restriction
\[
 g^2\mid T_{aP+i}\ \Longrightarrow\
 \deg g=u+vP,\quad0\le u\le\lfloor i/2\rfloor,
                       \quad0\le v\le\lfloor a/2\rfloor.       \tag{5}
\]
Every such $g$ has even degree: $T_n$ is the moment of
$(1+tX^2)^n$ under the positive Gamma shape $n$. Even $n$ has no real
root, and for odd $n$ the derivative integral is strictly positive on
$\mathbb R$. Hence all real roots are simple.
Residue square-factor degrees from (1) may further restrict $u,v$;
slopes and lengths alone do not retain that information.

\section{Irreducibility at every prime-power index}
\begin{theorem}\label{thm:irred}
For every prime $p$ and every $r\ge1$, $T_{p^r}$ is irreducible over
$\mathbb Q$ (in fact over $\mathbb Q_p$).
\end{theorem}
For $p\ge5$, Theorem~\ref{thm:shift} at $a=1,i=0$ has a single primitive
edge of width $P$, so every nonconstant factor must have degree $P$.
The two small primes require separate endpoint calculations.
\begin{proof}
For $p=2$, let $P=2^r$ and $1\le j\le P$. Write $v_2(j)=t$.
The exact valuation $v_2\binom Pj=r-t$ follows from
$\binom Pj=(P/j)\binom{P-1}{j-1}$ and the odd last binomial.
In the factorial block $(P)_{2j}$, the initial factor contributes $r$;
all subsequent factors $P+h$, $1\le h<2j$, have the valuation of $h$,
except $h=P$, which contributes one additional unit of valuation.
Thus
\[
 v_2(A_j)=2r-t+v_2((2j-1)!)+\mathbf1_{\{j>P/2\}}.
\]
At $j=P$ this is $2P-1$.
For $j<P$, binary digit sum gives
\[
 v_2(A_j)=2j+2r-1-s_2(j)-2t+\mathbf1_{\{j>P/2\}}
                                                    \ge2j,
\]
since $s_2(j)\le r-t$ and $t\le r-1$.
The polygon is therefore the single primitive edge from $(0,0)$ to
$(P,2P-1)$.

For $p=3$, the same binomial argument gives
$v_3\binom Pj=r-v_3(j)$. Now $v_3(P+h)=v_3(h)$ for every
$1\le h<2P$, including $h=P$, so
\[
                 v_3(A_j)=2r-v_3(j)+v_3((2j-1)!).
\]
At $j=P$ this is $P-1$.
For $j<P$, since $2j-1<2P$ has ternary digit sum at most $2r+1$,
\[
 v_3(A_j)=j+2r-v_3(j)-\frac{1+s_3(2j-1)}2\ge j.
\]
The polygon is the single primitive edge to $(P,P-1)$.
In both cases primitivity forces every nonconstant $\mathbb Q_p$ factor
to have full degree, proving the result.
\end{proof}

\begin{corollary}\label{cor:families}
Under the hypotheses of Theorem~\ref{thm:shift}, $T_{aP+i}$ is squarefree
for each of the following complete index families:
\[
 (a,i)\in\{(1,0),(1,1),(1,2),(1,3),
          (2,0),(2,1),(2,2),(3,0),(3,1),(4,0)\}.             \tag{6}
\]
The original cubic two moments at terminal $aP+i$ cannot both vanish.
\end{corollary}
In (5), $P$ is odd, so parity removes all positive patterns in (6)
except full-square patterns at $(a,i)=(2,2),(4,0)$.
Those patterns would make all of $T_n$ a rational square, since its
constant term is one. But its leading coefficient
$(3n-1)!/(n-1)!$ is not a square: Bertrand's theorem provides a prime
between $\lfloor(3n-1)/2\rfloor$ and twice that integer, occurring exactly
once in this ratio. Here $n\ge2$; the prime exceeds $n-1$.
This proves the corollary. It does not assert all multipliers or offsets.

\section{Exact analytic geometry: two nondegenerate folds}
With $t=w/n^2$, the integral is exactly
\[
 T_n(w/n^2)=\frac{n^n}{\Gamma(n)}
 \int_0^\infty s^{-1}\exp(n\Phi(s,w))\,ds,
 \qquad \Phi=\log s+\log(1+ws^2)-s.                       \tag{7}
\]
The logarithmic expression is used only in local nonzero charts;
the original integrand is polynomial times $s^{n-1}e^{-ns}$.
The saddle equation and its discriminant are
\[
 F(s,w)=ws^3-3ws^2+s-1=0,\qquad
 \operatorname{Disc}_s F=-4w(1-9w+27w^2).                 \tag{8}
\]
At $w=0$ the polynomial saddle equation drops degree and has the unique
finite saddle $s=1$; this is not an ordinary finite fold.
The finite collision points are
\[
 s_\pm=\frac{3\pm i\sqrt3}2,\qquad
 w_\pm=\frac16\pm\frac{i\sqrt3}{18},\qquad
 |w_\pm|=\frac1{3\sqrt3}.                                \tag{9}
\]
At either matching pair $(s_\pm,w_\pm)$,
\[
 \Phi_s=\Phi_{ss}=0,\quad
 \Phi_{sss}=-\frac{3(s_\pm-1)^2}{s_\pm^3}\ne0,
 \quad \Phi_{sw}=\frac{2s_\pm}{(1+w_\pm s_\pm^2)^2}\ne0,
                                                               \tag{10}
\]
with moduli $1/\sqrt3$ and $3\sqrt3/2$ respectively.
The third saddle is $3-2s_\pm=\mp i\sqrt3$.
Thus the exact phase has ordinary cubic folds, not separated simple
saddles everywhere. No contour-selection, Airy approximation, uniform
zero exclusion or effective simplicity cutoff is proved in this section.
A future analytic proof must justify all contributions and control every
root in the fold regions as well as away from them.

\section{Retained arithmetic limits and current scope}
At $n=732$, the slope-and-length square-factor degree sets for
$p=2,3,5,7,11$ still permit degree six; prime thirteen excludes it.
This disproves coverage by THAT fixed small set. It is not a failure of
all primes through the complete slope cap $3n-1$, or a counterexample
to SFC. The exact profiles are retained in the diagnostic file.
The tempting sufficient conditions $p>2k+1$, $n\bmod p<2k$, and even
$p>3k$, $n\bmod p<2k$, also fail as pure slope exclusions: degree four
survives at $(n,p)=(51,11)$ and $(118,13)$ respectively.
No actual repeated factor is inferred from either diagnostic.
A separate tempting shortcut also fails: the residual square kernel
$K_{11}$ from the preceding proof has degree five but exactly three real
roots, all positive, by an exact Sturm count. Its alternating coefficients
exclude negative roots. Thus universal real-rootedness of these kernels
cannot provide the missing simplicity theorem; this is not a repeated-root
or SFC counterexample.

For the integer coefficients of $T_n$, primes above $3n-1$ are all
units and yield only a horizontal polygon. Thus $3n-1$ is an exact
finite cap for slope-and-length restrictions at a supplied index;
residual tests above this cap can still supply information.
The diagnostic's finite replays are bookkeeping; the all-$r$ evidence
is the written valuation and digit argument, not extrapolation.
General $T_n$ simplicity, uniform prime/residual coverage, arbitrary
monomial supports, general FC and unrestricted SFC remain open.
Independent replay, formal verification, external review and priority
claims all remain false.
\end{document}
